\section{Conclusiones}

El proyecto dispone de una cadena reproducible desde una fuente raw preservada hasta un dataset
analítico de 89.740 canciones únicas. La preparación resolvió la anomalía con nulos y duración cero,
retiró el índice exportado, consolidó IDs y conservó géneros múltiples. Los controles posteriores
confirman el contrato esperado y las pruebas automatizadas protegen la lógica central.

El EDA demuestra que la popularidad se concentra en valores bajos y medios y que las asociaciones
acústicas individuales son débiles. \code{instrumentalness} ofrece la señal de mayor magnitud, pero
no determina el resultado. La redundancia entre energía, sonoridad y acousticness debe considerarse
en modelos lineales. Las diferencias por explicitud y género son descriptivas, presentan
solapamiento o sesgos de cobertura y no justifican afirmaciones causales.

La robustez adicional confirma que el efecto estandarizado de \code{explicit} es pequeño, que el
ranking principal de géneros no cambia entre soportes de 100, 200 y 500, y que artistas, álbumes y
nombres tienen cardinalidad suficiente para introducir memorización. Los extremos se mantienen
porque no existe evidencia general de error; se proponen transformaciones y validación de
sensibilidad.

Los límites más importantes son temporalidad desconocida, muestra no representativa, objetivo
dinámico y variables contextuales omitidas. Estos límites acotan tanto la capacidad predictiva como
la interpretación de negocio. Éticamente, una futura herramienta debe apoyar revisión humana,
explicitar incertidumbre y evitar automatizar asignaciones de alto impacto.

La siguiente etapa debe establecer una línea base de regresión, aplicar pipelines sin leakage,
codificar categorías según su semántica y comparar particiones aleatorias, agrupadas por artista y,
si se obtienen fechas, temporales. El éxito no consiste solo en reducir error global: requiere
generalización defendible, desempeño por segmentos, utilidad real y uso responsable.
